\section{Executive summary}
\label{sec:summary}

\textbf{Recommendation.} DataCo's late deliveries are mostly a promise-setting problem, not an
order-selection problem. We recommend three actions.
\begin{enumerate}
  \item \textbf{Change two promises first.} First Class is promised in 1 day but every First Class
        order takes exactly 2 days, so all of them are late. Same Day orders placed after 12:00
        always arrive the next day. Promising 2 days for First Class, and next-day delivery for
        afternoon Same Day orders, cuts late orders from \textbf{246 to 169 per week ($-31\%$)}.
        The average promise lengthens by only 0.18 days.
  \item \textbf{At order release, flag orders with a five-row lookup table, not a
        machine-learning model.} Flag an order when $P(\text{late})\times\text{order value} >
        \usd{200}$. This flags about 187 of 428 orders per week, catches 129 late ones and
        saves about \textbf{\usd{2{,}400} per week}, against \usd{1{,}500} for the best fixed rule.
  \item \textbf{Re-check on real delivery records before deploying.} The delivery timings in this
        dataset follow deterministic, apparently generated patterns.
\end{enumerate}

\textbf{Evidence.} We trained four baselines and six models (decision tree, random forest,
gradient boosting, SVM, MLP and logistic regression) on 49{,}826 orders from 2015--2017-06,
and tested them once on 13{,}071 later orders. The best models, the decision tree (PR-AUC
$0.854 \pm 0.005$ over 5 rolling-origin folds) and gradient boosting ($0.853 \pm 0.003$), tie
with the lookup table ($0.855 \pm 0.002$). Within Second and Standard Class the models rank
orders at chance level (ROC-AUC 0.47--0.52), and across 36 cost scenarios the best model and
the table stay within \usd{25} per week of each other. Shipping mode and order hour carry the
entire signal.
